%%=============================================================================
%% Surveyonderzoek
%%=============================================================================

\chapter{Surveyonderzoek}%
\label{ch:surveyanalyse}

Dit hoofdstuk rapporteert de resultaten van het kwantitatief surveyonderzoek dat
werd uitgevoerd ter aanvulling van de semigestructureerde interviews en de
vergelijkende app-analyse. De survey had een tweeledige doelstelling: enerzijds
het valideren van de kwalitatieve interviewbevindingen op een bredere steekproef,
anderzijds het verzamelen van directe gebruikersvoorkeuren over concrete
designkeuzes in het prototype.

\section{Opzet en verspreiding}

De survey werd opgesteld in Google Forms en bestond uit 23 vragen verdeeld over
vijf blokken: profielvragen, vragen over huidige planningstools, designkeuzes via
A/B-vergelijkingen, prototype-evaluatie op basis van Likert-schalen en een open
eindvraag. De volledige vragenlijst is opgenomen als bijlage~\ref{bijl:surveyanalyse}.

De survey werd verspreid via drie kanalen: de studentengemeenschap van HOGENT
(via Toledo, studentenkring en studentenbegeleiding), Reddit (\textit{r/belgium}
en \textit{r/ADHD}) en Discord. De oproep richtte zich uitdrukkelijk niet alleen
tot studenten met een officiële diagnose, maar ook tot personen die zich herkennen
in planningsproblemen of uitstelgedrag. In totaal namen \textbf{32 respondenten}
deel aan de survey.

\section{Profiel van de respondenten}

\subsection{ADHD-status}

Van de 32 respondenten had 56,3\% (n=18) een officiële ADHD-diagnose, 25,0\%
(n=8) vermoedt sterk ADHD te hebben zonder formele diagnose, en 18,8\% (n=6)
herkent zich in planningsproblemen zonder vermoeden van ADHD. Geen enkele
respondent gaf aan geen enkele planningsproblematiek te ervaren. De steekproef
is daarmee vrijwel volledig relevant voor het onderzoek.

\subsection{Relatie met deadlines}

Gevraagd naar hun relatie met deadlines, gaf 68,8\% (n=22) aan te weten dat ze
moeten beginnen maar er niet aan toe te komen, en 62,5\% (n=20) geeft aan pas
te starten als de deadline dichtbij is. Slechts 9,4\% (n=3) begint ruim op tijd
en werkt gespreid. Deze verdeling bevestigt het beeld dat naar voren kwam in de
interviews: het probleem ligt niet bij het kennen van de deadline, maar bij het
initiëren van de taak.

\subsection{Moeilijkste moment in het planningsproces}

Op de vraag wat het moeilijkste moment is in een planningsproces, antwoordde
77,4\% (n=24 van 31 respondenten) dat dit het starten is, ook als ze de wil
hebben om te beginnen. Dit is de sterkste kwantitatieve bevestiging van de
initiatieparalyse die in de literatuurstudie~\autocite{Barkley2022} en in beide
interviews naar voren kwam. Tijdsinschatting werd door 16,1\% als moeilijkst
aangeduid; verdergaan na een onderbreking en bijhouden wat er gedaan is, werden
elk door minder dan 5\% vermeld.

\subsection{Reactie op een gemiste deadline}

Op de vraag hoe men reageert wanneer een deadline gemist wordt, antwoordde
slechts 9,4\% (n=3) dat men meteen herplant en verdergaat. De overige 84,4\%
vertoont een of andere vorm van vermijding: 46,9\% doet de taak op het
allerlaatste moment alsnog, 37,5\% stelt herplannen uit door schaamte of stress,
en 6,3\% geeft de taak helemaal op.

Dit cijfer is de meest directe kwantitatieve onderbouwing van de kernlacune die
in de app-analyse werd geïdentificeerd: het moment ná een gemiste deadline is
het minst ondersteunde en tegelijk het meest kritieke interventievenster. Geen
enkele onderzochte planningsapp begeleidt actief dit moment, en de surveydata
toont aan dat 84,4\% van de gebruikers op dat moment ondersteuning nodig heeft.

\section{Huidige planningstools}

\subsection{Bekendheid en gebruik}

Google Calendar en de papieren agenda zijn de meest gebruikte tools, beiden door
75,0\% (n=24) van de respondenten. Microsoft To Do en Todoist werden elk door
34,4\% (n=11) geprobeerd, Notion door 28,1\% (n=9). Opmerkelijk is dat de
ADHD-specifieke apps Tiimo en Goblin.tools elk slechts door 3,1\% (n=1) van
de respondenten werden geprobeerd. Dit suggereert dat het bestaande ADHD-specifieke
aanbod de doelgroep niet of nauwelijks bereikt.

\subsection{Retentie}

50,0\% van de respondenten houdt een planningsapp minder dan vier weken vol:
25,0\% stopt binnen de eerste week en 25,0\% binnen één tot vier weken. Slechts
21,9\% volhoudt langer dan drie maanden. Dit retentieprobleem bevestigt de
bevinding uit de interviews dat apps structureel worden verlaten, niet sporadisch.

\subsection{Redenen van afhaken}

De voornaamste reden om met een planningsapp te stoppen is dat de app niet helpt
bij het daadwerkelijk starten van taken (35,5\%, n=11). Op de tweede plaats staat
het stoppen na een gemiste dag of deadline (25,8\%, n=8), gevolgd door te veel
moeite om informatie in te voeren (22,6\%, n=7) en notificaties die worden
genegeerd (16,1\%, n=5). Gecombineerd gaat 61,3\% van het afhaken terug op
vastlopen of op een gemiste deadline. Dit zijn twee situaties die het prototype expliciet
adresseert.

\subsection{Perceptie van bestaand aanbod}

Op de vraag in hoeverre bestaande planningsapps rekening houden met hoe het
brein van de respondent werkt (schaal 1--5), gaf 40,6\% een score van 1 en
37,5\% een score van 2. Niemand gaf een score van 4. Eén respondent gaf een
score van 5. Het gewogen gemiddelde bedraagt 1,88 op 5. Dit is de meest
uitgesproken score in de volledige survey en vormt een sterke kwantitatieve
aanwijzing dat het bestaande aanbod de doelgroep fundamenteel niet bereikt.

\section{Designkeuzes}

In het derde blok van de survey werden respondenten gevraagd te kiezen tussen
twee concrete ontwerpopties per vraag. De resultaten geven directe sturing aan
de ontwerpbeslissingen in het prototype.

\subsection{Taak invoeren}

Op de vraag hoe men een taak wil invoeren: via vrije tekstinvoer waarbij de app
daarna pas details vraagt (optie A), of via een gestructureerd formulier met
alle velden tegelijk (optie B) — koos 53,1\% voor optie B en 46,9\% voor optie
A. Het verschil is nipt. De implicatie voor het ontwerp is dat een combinatie
van beide benaderingen wenselijk is: een laagdrempelig tekstveldje als startpunt,
gevolgd door een overzichtelijke structuur.

\subsection{AI-taakverdeling}

De voorkeur voor hoe de AI een taak opsplitst is veel uitgesproken: 84,4\%
verkiest logische stappen op basis van wat samenhangt (bv. ``lezen'',
``samenvatten'', ``schrijven''), tegenover 15,6\% dat vaste tijdsblokken
verkiest. Dit bevestigt de ontwerpkeuze in het prototype om logische stappen
te genereren in plaats van tijdgestuurde blokken. Vaste tijdsblokken werden
als rigide en onrealistisch beschouwd.

\subsection{Reactie op een gemiste deadline}

62,5\% verkiest een neutrale herstart: de app toont de volgende ochtend één
rustige melding zonder cumulerende achterstanden. 37,5\% geeft de voorkeur
aan een volledig overzicht van alles wat gemist werd. Dit bevestigt de
ontwerpkeuze in het prototype om geen rode tellers of opeenstapelende
meldingen te tonen na een gemiste deadline.

\subsection{Herinnering aan taken}

71,9\% verkiest passieve zichtbaarheid via een widget op het startscherm boven
een actieve dagelijkse pushmelding (28,1\%). Dit sluit aan bij de bevinding van
respondent 2 in de interviews, die aangaf dat widgets beter werken dan
notificaties omdat ze zichtbaar zijn zonder actieve interactie te vereisen.
Voor een productieversie van de applicatie is widget-ondersteuning dan ook een
prioriteit.

\subsection{Dagweergave}

65,6\% verkiest een prioriteitenlijst zonder vaste tijden boven een tijdlijn op
basis van uur (34,4\%). Dit is een opmerkelijke bevinding: het huidige prototype
implementeert een tijdlijnstructuur met tijdslots. Op basis van dit surveyresultaat
wordt aanbevolen om in een volgende iteratie een hybride dagweergave te
implementeren waarbij tijden optioneel zijn en de prioriteitsvolgorde centraal
staat.

\subsection{Reactie bij vastlopen}

58,1\% geeft de voorkeur aan verdere opsplitsing van de huidige stap in
micro-stappen wanneer men vastloopt, tegenover 41,9\% dat een alternatieve taak
voorgesteld wil krijgen. Het prototype implementeert de opsplitsingsoptie correct.
De relatief hoge score voor de alternatieve taak suggereert dat een
``sla-deze-stap-over''-functie een waardevolle aanvulling zou zijn.

\subsection{Absolute dealbreakers}

Op de vraag wat een app absoluut niet mag doen, selecteerde 68,8\% (n=22) hoge
invoerdrempel, 62,5\% (n=20) bestraffing voor gemiste taken (visueel of via
tellers), en 40,6\% (n=13) directieve instructies die de gebruiker vertellen
wat hij of zij nu moet doen. Slechts 3,1\% (n=1) gaf aan dat niets van
bovenstaande hem stoort. Deze drie ontwerpverboden waren reeds verankerd in de
vereisten van het prototype; de surveydata bevestigt dat ze correct
geïdentificeerd werden.

\section{Prototype-evaluatie}

\subsection{Gebruiksgemak}

30 respondenten testten het prototype. Op de uitspraak ``Ik vond het prototype
makkelijk te gebruiken'' (schaal 1--5) antwoordde 43,3\% neutraal (score 3) en
26,7\% positief (score 4). Het gewogen gemiddelde bedraagt 2,83 op 5. Dit is
een matig resultaat, maar realistisch voor een klikbaar HTML-prototype zonder
echte AI-integratie, mobiele optimalisatie of persistente dataopslag.

\subsection{Rust en overzichtelijkheid}

Op de uitspraak ``De interface voelde rustig en overzichtelijk aan'' scoorde de
applicatie een gewogen gemiddelde van 3,03 op 5, met 33,3\% neutraal en 33,3\%
positief. 16,7\% gaf een score van 1, wat aansluit bij de open antwoorden
waarin meerdere respondenten het kleurgebruik als overweldigend omschreven.

\subsection{Gepastheid van de herplanningsaanpak}

De uitspraak ``De manier waarop de app reageert op een gemiste deadline voelt
gepast aan voor mij'' behaalde eveneens een gemiddelde van 3,03 op 5. 33,3\%
scoorde positief (score 4) en 33,3\% neutraal (score 3). Dit suggereert dat de
niet-bestraffende herplanningsflow door een meerderheid als passend wordt
ervaren, al is er ruimte voor verbetering.

\subsection{Intentie tot gebruik}

Op de uitspraak ``Ik zou deze app willen gebruiken als hij beschikbaar was op
mijn telefoon'' antwoordde 51,7\% positief en 48,3\% negatief (n=29). De
meest waardevolle functies waren de dagweergave met kleurgecodeerde taken (30,0\%)
en het opsplitsen van taken in stappen via het taakdetailscherm (30,0\%),
gevolgd door de patroonherkenning (20,0\%).

\subsection{Grootste drempel voor dagelijks gebruik}

44,8\% (n=13) geeft aan de gewoonte te zullen verliezen om de app te openen als
grootste drempel voor dagelijks gebruik. 24,1\% mist een echte mobiele app in
plaats van een browserbestand, en 13,8\% vindt de invoer nog te omslachtig.
Dit retentieprobleem ligt niet bij de interface of de functies, maar bij de
afwezigheid van passieve zichtbaarheid, wat opnieuw pleit voor
widget-integratie als prioriteit in een productieversie.

\section{Open antwoorden}

14 respondenten vulden de open eindvraag in. De antwoorden werden thematisch
gecodeerd en leverden vijf thema's op.

\textbf{Agenda-integratie.} Meerdere respondenten vroegen om koppeling met
bestaande kalenders zoals Google Calendar. Dit is de meest genoemde ontbrekende
functionaliteit.

\textbf{Kleurgebruik en woordkeuze.} Het kleurgebruik werd door meerdere
respondenten als te intensief of als waardeoordeel ervaren. Eén respondent
omschreef de eerste indruk als ``een kleurboek''. Een andere respondent merkte
op dat de term ``kleiner maken'' kleinerend aanvoelt en stelde ``heropdelen''
voor als neutraler alternatief. Dit zijn bruikbare aanwijzingen voor een volgende
iteratie.

\textbf{Keuzereductie.} Meerdere respondenten vroegen om minder keuzes en
meer automatisering. Eén respondent beschreef het gewenste mechanisme expliciet:
de gebruiker geeft aan wanneer hij vrij is, en de app bepaalt autonoom wanneer
taken ingepland worden.

\textbf{Gamificatie.} Eén respondent pleitte voor gamificatie als
motivatiehefboom, met als argument dat mensen met ADHD vatbaar zijn voor
positieve bekrachtigingsmechanismen. Dit gaat in tegen de bewuste ontwerpkeuze
om punten en streaks te vermijden, en verdient verdere overweging in een
volgende versie.

\textbf{Grenzen van planningsapps.} Twee respondenten wezen erop dat
planningsapps het kernprobleem van executive dysfunction niet oplossen. Eén
respondent beschreef concrete triggers die een mentale blokkade veroorzaken —
onzekerheid, ontbrekende informatie, openstaande sociale verplichtingen — en
stelde een AI-chatbot voor die helpt de oorzaak van de blokkade te analyseren.
Dit is een eerlijke begrenzing van wat een planningsapp kan bieden, en een
richting voor toekomstig onderzoek.

\section{Besluit}

De surveyresultaten bevestigen en verdiepen de bevindingen uit de
literatuurstudie, de app-analyse en de interviews op drie niveaus.

Op het niveau van de \textbf{probleemstelling} wordt de kernlacune kwantitatief
onderbouwd: 84,4\% van de respondenten vertoont vermijding na een gemiste
deadline, slechts 9,4\% herplant meteen. 78,1\% scoort bestaande apps een 1
of 2 op 5 voor hun aansluiting bij de werking van het ADHD-brein.

Op het niveau van de \textbf{ontwerpvalidatie} bevestigen de A/B-resultaten de
meeste gemaakte keuzes: logische stappen (84,4\%), neutrale herstart bij gemiste
deadline (62,5\%), widget boven pushmelding (71,9\%) en geen bestraffing (62,5\%).
Één keuze wordt bijgestuurd: de dagweergave dient in een volgende iteratie een
prioriteitenlijst als primaire weergave te gebruiken in plaats van een tijdlijn,
op basis van de voorkeur van 65,6\% van de respondenten.

Op het niveau van de \textbf{prototypekwaliteit} liggen de gebruiksgemak- en
overzichtelijkheidsscores rond 3 op 5, wat realistisch is voor een klikbaar
HTML-prototype. De voornaamste verbeterpunten zijn het kleurgebruik, de
invoerdrempel en de afwezigheid van widget-ondersteuning. 51,7\% van de
respondenten geeft aan de app te willen gebruiken als hij beschikbaar was op
hun telefoon, wat voldoende intentie aantoont om verdere ontwikkeling te
rechtvaardigen.